\begin{recipe}
[%
    preparationtime = {\unit[30]{min}},
    portion = {\portion{4}},
    source = {\protect\recipesource{https://www.habe-ich-selbstgemacht.de/zitronenrisotto-rezept/}{Selbstgemacht: Zitronenrisotto}},
    calory = {694kcal | 18g Protein | 89g KH | 27g Fett},
    price = {1,70€},
    created = {24.06.2026},
    %rating = {1},
]
{Zitronenrisotto}

\graph{% pictures
    big=pic/hauptgericht/20_zitronenrisotto
}

\introduction{%
Ein cremiges Risotto mit frischer Zitrone, Weißwein und Parmesan. Die Mischung aus Säure, Butter und Käse macht es gleichzeitig leicht, intensiv und richtig schön schlotzig.
}

\ingredients{
    1 & Mittelgroße Zwiebel\\
    1 Zehe & Knoblauch\index{Gemüse!Knoblauch}\\
    2 EL & Olivenöl\\
    400 g & Risottoreis\index{Getreide!Risottoreis}\\
    130 ml & Trockener Weißwein\\
    1 l & Gemüsebrühe\\
    1 & Bio-Zitrone\index{Obst!Zitrone}\\
    60 g & Butter\\
    100 g & Parmesan, frisch gerieben\index{Milchprodukte \& Alternativen!Parmesan}\\
    & Salz, schwarzer Pfeffer
}

\preparation{%
    \step%
    Zwiebel und Knoblauch fein würfeln. Die Zitronenschale abreiben, den Saft auspressen und den Parmesan fein reiben. Gemüsebrühe in einem zweiten Topf erhitzen und warm halten.

    \step%
    Olivenöl in einem großen Topf erhitzen. Zwiebel und Knoblauch darin bei mittlerer Hitze glasig dünsten.

    \step%
    Risottoreis dazugeben und etwa 2 Minuten unter Rühren anschwitzen, bis die Körner am Rand leicht glasig werden.

    \step%
    Mit Weißwein ablöschen und fast vollständig einkochen lassen.

    \step%
    Nach und nach jeweils eine Kelle heiße Gemüsebrühe einrühren. Die nächste Portion erst zugeben, wenn der Reis die Flüssigkeit weitgehend aufgenommen hat. So etwa 20--25 Minuten weitergaren, bis der Reis cremig, aber noch leicht bissfest ist.

    \step%
    Topf vom Herd nehmen. Zitronenabrieb, Zitronensaft, Butter und Parmesan gründlich unterrühren. Mit Salz und schwarzem Pfeffer abschmecken und sofort servieren.
}

\hint{
    Die Brühe sollte beim Einrühren heiß sein, damit der Garprozess nicht immer wieder unterbrochen wird. Das fertige Risotto direkt servieren, solange es noch cremig fließt.
}

\suggestion[Topping wie auf dem Bild]{%
    Zucchini in Scheiben kräftig in Olivenöl anbraten, würzen und zusammen mit frischem Rucola auf dem Risotto anrichten.
}

\end{recipe}
